
\medskip
\phantomsection\label{par:sga3-VIA-4.4}%
\noindent\textbf{4.4.---}
It remains to construct the cokernel of the groupoid
\(\mathscr S(G^\alpha_*)\) when \(G^\alpha\) is an arbitrary
connected component of \(G\). Let \(A'\) be a sufficiently large
local \(A\)-algebra, finite and free over \(A\) (cf.\ \SGARefLink{thm:sga3-VIA-3.2}{3.2}). Then
\(G^\alpha\otimes_AA'\) is the union of finitely many connected
components
\[
  C^1,\ldots,C^n
\]
of \(G\otimes_AA'\), each of which has a strictly rational point.
For every \(i\), there is therefore a right translation \(r_i\) of
\(G\otimes_AA'\) which maps \(G^0\otimes_AA'\) onto \(C^i\).
This translation induces an isomorphism
\[
  \mathscr S(G^0_*)\otimes_AA'
    \xrightarrow{\ \sim\ }
  \mathscr S(C^i_*).
\]
Consequently, the groupoid induced by \(G_*\otimes_AA'\) on the
saturation of \(C^i\) has a cokernel.

Since \(\mathscr S(G^\alpha_*)\otimes_AA'\) is the direct sum of
some of the groupoids \(\mathscr S(C^i_*)\), it too has a cokernel.
This cokernel is a direct sum of some number of copies of
\[
  (F_0\otimes_AA')\backslash(G^0\otimes_AA').
\]
Thus every finite subset of it is contained in an affine open subset.
Moreover, formation of this cokernel commutes with finite locally
free extensions of the base (cf.\ editor's
note~\ref{note:sga3-VIA-37} and \SGARefLink{par:sga3-VIA-4.6}{4.6} below). As in \SGARefLink{par:sga3-VIA-3.2.3}{3.2.3}, it follows
that this cokernel has the form \(Y\otimes_AA'\), where \(Y\) is a
cokernel of \(\mathscr S(G^\alpha_*)\).
